% --- Вставити у §5.2 "Push-режим (фонове оновлення)" після вступного абзацу ---

Частота оновлення визначається як функція трьох сигналів активності:

\begin{equation}
  f_{\text{refresh}}(\text{task}) = g\bigl(
    \text{commits-touching-task},\;
    \text{comments-on-issue},\;
    \text{time-since-last-pull}
  \bigr)
  \label{eq:refresh}
\end{equation}

\noindent де $g$ --- монотонно зростаюча функція за кожним аргументом, обмежена зверху одним оновленням на 24~години та знизу --- одним оновленням на 7~днів для будь"=якої задачі з позначкою \texttt{LONG-TERM}.
Таке обмеження запобігає як надмірному витрачанню обчислювальних ресурсів на високоактивних задачах, так і повному застаріванню контексту для задач з низькою активністю.

Кожне оновлення породжує три артефакти:

\begin{itemize}[nosep]
  \item \textbf{Виконавче резюме змін} --- генерується мовною моделлю на основі дифів релевантних комітів та коментарів до задачі; містить стислий виклад усіх суттєвих подій від попереднього оновлення.
  \item \textbf{Дельта лінеажу інструментів} --- перелік нових, модифікованих та виведених з експлуатації інструментів із дифами параметричних схем, що дозволяє оператору при поверненні одразу бачити зміни в доступному арсеналі засобів.
  \item \textbf{Відкриті питання} --- нерозв'язані альтернативи та точки прийняття рішень, позначені сумаризатором як такі, що потребують втручання оператора при відновленні роботи над задачею.
\end{itemize}

За військовою аналогією, цей механізм відповідає автоматизованій підготовці брифінгу передачі зміни: система формує зведення навіть тоді, коли жоден оператор не перебуває на бойовому чергуванні, --- щоб наступна зміна при заступанні отримала актуальну обстановку без затримки на повторний збір інформації.
